\documentclass[12pt]{article}
\usepackage[margin=0.8in]{geometry}
\usepackage[spanish]{babel}
\usepackage[utf8]{inputenc}
\usepackage{setspace}
\usepackage{graphicx}
\usepackage{listings}
\usepackage{xcolor}
\usepackage{fancyhdr}

% Configuración de código
\lstset{
    language=Java,
    basicstyle=\ttfamily\small,
    keywordstyle=\color{blue},
    commentstyle=\color{gray},
    stringstyle=\color{red},
    breaklines=true,
    showstringspaces=false,
    tabsize=2,
    numbers=left,
    numberstyle=\tiny,
    frame=single,
    backgroundcolor=\color{white}
}

\title{\textbf{Patrones de Diseño en Space Invaders} \\
\large{Análisis detallado de Factory, Strategy y Composite}}
\author{}
\date{}

\begin{document}

\maketitle
\onehalfspacing
\newpage
\tableofcontents
\newpage

\section{Introducción}

En el desarrollo del proyecto \textit{Space Invaders} se han aplicado varios patrones de diseño fundamentales de la programación orientada a objetos. Estos patrones permiten estructurar el código de forma más clara, mantenible y extensible, siguiendo principios como SOLID (Single Responsibility, Open/Closed, Liskov Substitution, Interface Segregation, Dependency Inversion).

\subsection{¿Por qué utilizar patrones de diseño?}

Los patrones de diseño son soluciones probadas a problemas comunes en programación. En este proyecto, abordan tres desafíos principais:

\begin{enumerate}
    \item \textbf{Creación de objetos:} ¿Cómo generar múltiples tipos de naves de forma controlada y centralizada?
    \item \textbf{Variación de comportamientos:} ¿Cómo permitir que el jugador cambie entre diferentes tipos de disparos dinámicamente?
    \item \textbf{Composición de estructuras:} ¿Cómo representar de forma uniforme naves y disparos como colecciones de píxeles?
\end{enumerate}

En particular, se han utilizado tres patrones principales: \textbf{Factory}, \textbf{Strategy} y \textbf{Composite}, cada uno resolviendo un problema específico dentro del sistema.

\subsection{Relación entre patrones}

Los patrones no existen de forma aislada sino que interactúan entre sí:

\begin{itemize}
    \item \textbf{Factory} crea instancias de \texttt{Naves} que contienen \textbf{Composites}.
    \item \textbf{Strategy} define el comportamiento de los disparos que son envueltos en \textbf{Composites}.
    \item \textbf{Composite} proporciona la estructura unificada para representar tanto naves como disparos.
\end{itemize}

\newpage
\section{Patrón Factory: Creación Centralizada de Naves}

\subsection{Descripción General}

El patrón \textbf{Factory} es un patrón de creación cuyo objetivo principal es delegar la responsabilidad de instanciar objetos a una clase especializada. En lugar de crear objetos directamente mediante el operador \texttt{new}, el cliente solicita dichos objetos a una fábrica, que encapsula toda la lógica de instanciación.

\textbf{Problema que resuelve:}
\begin{itemize}
    \item Eliminación del acoplamiento entre el cliente y las clases concretas
    \item Centralización de parámetros de inicialización
    \item Facilidad para extender el sistema sin modificar código existente
    \item Control único sobre la creación de objetos
\end{itemize}

\subsection{Implementación en el Proyecto}

En \textit{Space Invaders}, el patrón Factory se implementa a través de la clase \texttt{NaveFactory}, que centraliza la creación de todas las naves jugables (\texttt{Nave1}, \texttt{Nave2}, \texttt{Nave3}).

Además, la clase implementa el patrón \textbf{Singleton}, garantizando que exista una única instancia de la fábrica durante toda la ejecución del programa.

\subsection{Diagrama de Clases}

\begin{center}
    \includegraphics[width=0.9\textwidth]{images/factory_diagram.png}
    
    \textit{Figura 1: Estructura de clases del patrón Factory. NaveFactory actúa como punto único de creación para Nave1, Nave2 y Nave3.}
\end{center}

\subsection{Estructura del Código}

\begin{lstlisting}
public class NaveFactory {
    private static NaveFactory myFactory;
    
    private NaveFactory() {}
    
    public static NaveFactory getNaveFactory() {
        if (myFactory == null) {
            myFactory = new NaveFactory();
        }
        return myFactory;
    }
    
    public Naves generate(String pNave) {
        int x = 50;           // Posición inicial
        int y = 55;           // Posición inicial
        int velocidad = 1;    // Velocidad
        
        if (pNave == null) {
            return new Nave1(x, y, velocidad);
        }
        
        switch (pNave) {
            case "Nave1":
                return new Nave1(x, y, velocidad);
            case "Nave2":
                return new Nave2(x, y, velocidad);
            case "Nave3":
                return new Nave3(x, y, velocidad);
            default:
                return null;
        }
    }
}
\end{lstlisting}

\subsubsection{Explicación de los componentes}

\begin{description}
    \item[\texttt{myFactory}] Variable estática que almacena la única instancia de la fábrica.
    
    \item[\textit{Constructor privado}] Impide que el cliente instancie directamente \texttt{NaveFactory}. Solo se puede crear a través de \texttt{getNaveFactory()}.
    
    \item[\texttt{getNaveFactory()}] Implementa el patrón Singleton. Verifica si la instancia existe; si no, la crea. En llamadas posteriores, retorna la instancia exitente.
    
    \item[\texttt{generate(String)}] Método factory que recibe un identificador y retorna la nave correspondiente. Parámetros iniciales sde inicialización (\texttt{x}, \texttt{y}, \texttt{velocidad}) se definen aquí de forma centralizada.
\end{description}

\subsection{Flujo de Ejecución}

El flujo de creación de una nave sigue estos pasos:

\begin{enumerate}
    \item \textbf{Cliente solicita nave:}
    \begin{lstlisting}
NaveFactory fabrica = NaveFactory.getNaveFactory();
Naves nave = fabrica.generate("Nave1");
    \end{lstlisting}
    
    \item \textbf{Verificación Singleton:} \texttt{getNaveFactory()} comprueba si \texttt{myFactory} es nulo. La primera vez lo crea; después devuelve la existente.
    
    \item \textbf{Identificación de tipo:} El método \texttt{generate()} recibe el string \texttt{"Nave1"}.
    
    \item \textbf{Inicialización de parámetros:} Se asignan valores por defecto: posición (50, 55) y velocidad 1.
    
    \item \textbf{Instanciación:} Se crea una nueva instancia de \texttt{Nave1} con esos parámetros.
    
    \item \textbf{Retorno:} Se devuelve como tipo \texttt{Naves} (clase abstracta padre), permitiendo polimorfismo.
    
    \item \textbf{Inicialización de estructuras:} El constructor de \texttt{Nave1} llama a \texttt{construir()}, que crea su estructura de píxeles (patrón Composite).
\end{enumerate}

\textbf{Diagrama de secuencia conceptual:}

\begin{verbatim}
Cliente
   |
   |-- getNaveFactory() ---> Instancia Singleton
   |
   |-- generate("Nave1") -----> Verifica en switch
   |
   |-- new Nave1(50, 55, 1) ---|
   |                            |
   |                            v
   |                        Constructor Nave1
   |                            |
   |                            |-- inicializarNaveJugador()
   |                            |-- construir()
   |                            |
   |<-- Retorna Naves ---------|
   |
   v
(Nave lista para usar)
\end{verbatim}

\subsection{Ventajas del Patrón Factory}

\begin{enumerate}
    \item \textbf{Encapsulación:} La lógica de creación queda centralizada. Cambios en los parámetros iniciales se hacen en un único lugar.
    
    \item \textbf{Bajo acoplamiento:} El cliente solo depende de la interfaz \texttt{Naves}, no de las clases concretas.
    
    \item \textbf{Extensibilidad:} Agregar una nueva nave (\texttt{Nave4}) es sencillo:
    \begin{lstlisting}
// Solo añadir dos líneas:
case "Nave4":
    return new Nave4(x, y, velocidad);
    \end{lstlisting}
    
    \item \textbf{Singleton integrado:} Garantiza una única instancia de fábrica durante toda la ejecución.
    
    \item \textbf{Control global:} Si en el futuro se necesitan parámetros dinámicos (dificultad aumenta velocidad), todo se modifica en un único lugar.
\end{enumerate}

\subsection{Antipatrones evitados}

\textbf{Sin Factory (Acoplado):}
\begin{lstlisting}
// En MainFrame o Main
if (tipoNave.equals("Nave1")) {
    nave = new Nave1(50, 55, 1);
} else if (tipoNave.equals("Nave2")) {
    nave = new Nave2(50, 55, 1);
} else if (tipoNave.equals("Nave3")) {
    nave = new Nave3(50, 55, 1);
}
// Duplicación de parámetros en múltiples lugares
// Difícil de mantener
\end{lstlisting}

\textbf{Con Factory (Limpio):}
\begin{lstlisting}
Naves nave = NaveFactory.getNaveFactory()
    .generate(tipoNave);
// Una única línea, lógica centralizada
\end{lstlisting}

\newpage
\section{Patrón Strategy: Comportamiento de Disparos Intercambiable}

\subsection{Descripción General}

El patrón \textbf{Strategy} es un patrón de comportamiento que permite definir una familia de algoritmos, encapsularlos en clases independientes y hacerlos intercambiables en tiempo de ejecución. El cliente puede cambiar de estrategia sin conocer los detalles de implementación.

\textbf{Problema que resuelve:}
\begin{itemize}
    \item Múltiples formas de realizar la misma tarea (diferentes tipos de disparos)
    \item Necesidad de cambiar comportamiento dinámicamente sin recompilar
    \item Eliminar código condicional complejo (múltiples \texttt{if-else})
    \item Gestión independiente del estado de cada estrategia (munición)
\end{itemize}

\subsection{Implementación en el Proyecto}

En \textit{Space Invaders}, el patrón Strategy se emplea para gestionar los distintos tipos de disparo disponibles. Se define una interfaz común \texttt{StrategyDisparo} que todas las estrategias deben implementar.

\subsection{Diagrama de Clases}

\begin{center}
    \includegraphics[width=0.85\textwidth]{images/strategy_diagram.png}
    
    \textit{Figura 2: Estructura de Strategy. StrategyDisparo define el contrato; cada implementación maneja su propia munición.}
\end{center}

\subsection{Interfaz Strategy}

\begin{lstlisting}
public interface StrategyDisparo {
    // Obtiene identificador del tipo de disparo
    String getTipo();
    
    // Obtiene munición disponible (-1 = infinita)
    int getMunicion();
    
    // Consume una unidad de munición
    // Retorna false si no queda
    boolean gastar();
    
    // Verifica si hay munición disponible
    boolean tieneMunicion();
}
\end{lstlisting}

\subsection{Implementaciones Concretas}

\subsubsection{DisparoPixel - Munición Infinita}

\begin{lstlisting}
public class DisparoPixel implements StrategyDisparo {
    @Override
    public String getTipo() {
        return "pixel";
    }
    
    @Override
    public int getMunicion() {
        return -1;  // Infinita
    }
    
    @Override
    public boolean gastar() {
        return true;  // Siempre disponible
    }
    
    @Override
    public boolean tieneMunicion() {
        return true;  // Nunca se agota
    }
}
\end{lstlisting}

\textbf{Características:}
\begin{itemize}
    \item Tipo: \texttt{"pixel"}
    \item Munición: Infinita (-1)
    \item Disponibilidad: Siempre puede disparar
    \item Uso: Disparo base predeterminado
\end{itemize}

\subsubsection{DisparoFlecha - Munición Limitada (30 unidades)}

\begin{lstlisting}
public class DisparoFlecha implements StrategyDisparo {
    private int municion = 30;
    
    @Override
    public String getTipo() {
        return "flecha";
    }
    
    @Override
    public int getMunicion() {
        return municion;
    }
    
    @Override
    public boolean gastar() {
        if (municion > 0) {
            municion--;
            return true;
        }
        return false;  // No hay más munición
    }
    
    @Override
    public boolean tieneMunicion() {
        return municion > 0;
    }
}
\end{lstlisting}

\textbf{Características:}
\begin{itemize}
    \item Tipo: \texttt{"flecha"}
    \item Munición inicial: 30 unidades
    \item Comportamiento: Decrementa munición cada vez que se dispara
    \item Disponibilidad: Solo mientras haya munición
\end{itemize}

\subsubsection{DisparoRombo - Munición Limitada (20 unidades)}

\begin{lstlisting}
public class DisparoRombo implements StrategyDisparo {
    private int municion = 20;
    
    @Override
    public String getTipo() {
        return "rombo";
    }
    
    @Override
    public int getMunicion() {
        return municion;
    }
    
    @Override
    public boolean gastar() {
        if (municion > 0) {
            municion--;
            return true;
        }
        return false;
    }
    
    @Override
    public boolean tieneMunicion() {
        return municion > 0;
    }
}
\end{lstlisting}

\textbf{Características:}
\begin{itemize}
    \item Tipo: \texttt{"rombo"}
    \item Munición inicial: 20 unidades (menos que flecha)
    \item Comportamiento: Similar a flecha pero con recursos más limitados
    \item Disponibilidad: Limitada a 20 disparos
\end{itemize}

\subsection{Uso en Naves}

Cada nave posee sus estrategias de disparo permitidas:

\begin{lstlisting}
public class Nave1 extends Naves {
    private DisparoPixel estrategiaPixel = 
        new DisparoPixel();
    private DisparoFlecha estrategiaFlecha = 
        new DisparoFlecha();
    
    @Override
    public ArrayList<StrategyDisparo> 
            getEstrategiasPermitidas() {
        ArrayList<StrategyDisparo> estrategias = 
            new ArrayList<>();
        estrategias.add(estrategiaPixel);
        estrategias.add(estrategiaFlecha);
        return estrategias;
    }
}
\end{lstlisting}

\subsection{Flujo de Ejecución: Disparar con Strategy}

El proceso de disparar utiliza completamente el patrón Strategy:

\begin{enumerate}
    \item \textbf{Jugador selecciona tipo de disparo} (por ej., "flecha")
    
    \item \textbf{Sistema obtiene la estrategia:}
    \begin{lstlisting}
StrategyDisparo estrategiaActual = 
    nave.getEstrategia("flecha");
    \end{lstlisting}
    
    \item \textbf{Verificación de munición:}
    \begin{lstlisting}
if (estrategiaActual.tieneMunicion()) {
    // Paso 4
} else {
    // No hay munición, no disparar
}
    \end{lstlisting}
    
    \item \textbf{Consumir munición:}
    \begin{lstlisting}
estrategiaActual.gastar();
    \end{lstlisting}
    
    \item \textbf{Crear proyectil con forma según estrategia:}
    \begin{lstlisting}
String tipo = estrategiaActual.getTipo();
// "flecha" -> 3 píxeles en forma de flecha
// "rombo" -> 5 píxeles en forma de rombo
// "pixel" -> 1 píxel simple
Disparo nuevoDisparo = crearDisparo(tipo);
    \end{lstlisting}
    
    \item \textbf{Adicionar a lista de disparos:}
    \begin{lstlisting}
disparos.add(nuevoDisparo);
    \end{lstlisting}
    
    \item \textbf{Siguiente frame:} El disparo se mueve y se detectan colisiones.
\end{enumerate}

\textbf{Diagrama de secuencia conceptual:}

\begin{verbatim}
Judagador presiona "disparar" + "flecha"
   |
   v
Juego verifica estrategia actual
   |
   v
¿Hay munición? -----NO---> Mostrar "Sin munición"
   |
   SI
   |
   v
estrategiaActual.gastar() -----> municion = 30 a 29
   |
   v
Crear cuerpo de disparo
   (forma según getTipo())
   |
   v
Ejecutar update() cada frame
   |
   v
Disparos suben (mover -Y)
   |
   v
Detectar colisión
   |
   v
Eliminar disparo
\end{verbatim}

\subsection{Ventajas del Patrón Strategy}

\begin{enumerate}
    \item \textbf{Intercambiabilidad en tiempo de ejecución:} El jugador puede cambiar entre disparos sin pausar ni recargar.
    
    \item \textbf{Encapsulación de comportamientos:} Cada estrategia gestiona su propia lógica (munición, forma, comportamiento).
    
    \item \textbf{Reducción de condicionales:} Sin Strategy habría complejos \texttt{if-else} evaluándose constantemente:
    \begin{lstlisting}
// SIN STRATEGY (malo)
if (tipoDisparo == 1) {
    // Lógica pixel...
} else if (tipoDisparo == 2) {
    // Lógica flecha...
} else if (tipoDisparo == 3) {
    // Lógica rombo...
}
    \end{lstlisting}
    
    \item \textbf{Extensibilidad:} Nueva estrategia (ej. \texttt{DisparoLáser}) solo requiere implementar la interfaz.
    
    \item \textbf{Testing:} Cada estrategia se prueba independientemente.
\end{enumerate}

\newpage
\section{Patrón Composite: Estructura Jerárquica de Píxeles}

\subsection{Descripción General}

El patrón \textbf{Composite} es un patrón estructural que permite componer objetos en estructuras jerárquicas en forma de árbol. Su principal ventaja es que permite tratar de forma uniforme tanto a los objetos individuales como a las composiciones de objetos.

\textbf{Problema que resuelve:}
\begin{itemize}
    \item Representación de estructuras jerárquicas (naves, disparos complejos)
    \item Operaciones uniformes en elementos simples o complejos
    \item Movimiento coordiado de múltiples píxeles
    \item Detección de colisiones centralizada
    \item Flexibilidad geométrica para crear formas variadas
\end{itemize}

\subsection{Implementación en el Proyecto}

En este proyecto, el patrón Composite se utiliza para representar la estructura de naves y disparos como árboles de píxeles. Se define una interfaz común \texttt{Component} implementada por píxeles individuales (\texttt{Pixel}) y agrupaciones (\texttt{Composite}).

\subsection{Diagrama General}

\begin{center}
    \includegraphics[width=1\textwidth]{images/composite_structure.png}
    
    \textit{Figura 3: Estructura general del proyecto mostrando cómo Composite y Component integran con Disparo y Naves.}
\end{center}

\subsection{Interfaz Component}

\begin{lstlisting}
public interface Component {
    // Desplaza el componente dx píxeles en X, dy en Y
    void mover(int dx, int dy);
    
    // Obtiene coordenada X de referencia
    int getRefX();
    
    // Obtiene coordenada Y de referencia
    int getRefY();
    
    // Verifica si el componente está activo
    // (inactivo = fuera del tablero o destruido)
    boolean isActivo();
    
    // Activar/desactivar componente
    void setActivo(boolean b);
    
    // Métodos default que notifican evento al tablero
    default void notificarMovimientoJugador(int oldX, int oldY, 
                                             int newX, int newY) {
        Espacio.getEspacio().notificarMovimientoJugador
            (oldX, oldY, newX, newY);
    }
    
    default void notificarMovimientoDisparo(int oldX, int oldY, 
                                            int newX, int newY) {
        Espacio.getEspacio().notificarMovimientoDisparo
            (oldX, oldY, newX, newY);
    }
    
    default void notificarDisparoNuevo() { }
}
\end{lstlisting}

\subsection{Clase Hoja: Pixel}

\texttt{Pixel} representa una celda individual del tablero. Es la unidad atómica del Composite:

\begin{lstlisting}
public class Pixel implements Component {
    private int x;
    private int y;
    private boolean activo = true;
    private final boolean proyectil;
    
    // Constructor para píxeles de nave
    public Pixel(int x, int y) {
        this(x, y, false);
    }
    
    // Constructor con flag diferenciador
    public Pixel(int x, int y, boolean proyectil) {
        this.x = x;
        this.y = y;
        this.proyectil = proyectil;
    }
    
    @Override
    public void mover(int dx, int dy) {
        if (proyectil) {
            if (!activo) return;
            int oldX = x;
            int oldY = y;
            x += dx;
            y += dy;
            
            // Desactivar si sale del tablero
            if (y < 0) activo = false;
            
            // Notificar movimiento
            notificarMovimientoDisparo(oldX, oldY, x, y);
        } else {
            // Píxel de nave: solo actualizar posición
            x += dx;
            y += dy;
        }
    }
    
    @Override
    public int getRefX() { return x; }
    
    @Override
    public int getRefY() { return y; }
    
    @Override
    public boolean isActivo() { return activo; }
    
    @Override
    public void setActivo(boolean b) { this.activo = b; }
    
    @Override
    public void notificarDisparoNuevo() {
        if (proyectil) {
            Espacio.getEspacio().notificarDisparoNuevo(
                getRefX(), getRefY());
        }
    }
}
\end{lstlisting}

\textbf{Características:}
\begin{itemize}
    \item \textbf{Bandera \texttt{proyectil}:} Distingue píxeles de nave vs disparo (comportamientos diferentes).
    \item \textbf{Píxeles de nave:} Solo actualizan posición, no notifican movimiento individual.
    \item \textbf{Píxeles de disparo:} Notifican cada movimiento al tablero para detección de colisiones.
    \item \textbf{Desactivación automática:} Píxeles de disparo se desactivan al superar el límite superior del tablero.
\end{itemize}

\subsection{Clase Compuesta: Composite}

\texttt{Composite} agrupa múltiples \texttt{Components} (Pixels u otros Composites) en una estructura jerárquica:

\begin{lstlisting}
public class Composite implements Component {
    private final List<Component> components = 
        new ArrayList<>();
    private final boolean proyectil;
    
    // Constructor para Composite de nave (raíz)
    public Composite() {
        this(false);
    }
    
    // Constructor con diferenciador
    public Composite(boolean proyectil) {
        this.proyectil = proyectil;
    }
    
    public void addComponent(Component c) {
        components.add(c);
    }
    
    public void removeComponent(Component c) {
        components.remove(c);
    }
    
    public List<Component> getComponents() {
        return components;
    }
    
    // Obtiene todas las celdas ocupadas por píxeles activos
    // Recursivo: maneja agrupamientos anidados
    public List<int[]> celdasOcupadasActivas() {
        List<int[]> celdas = new ArrayList<>();
        for (Component c : components) {
            if (c instanceof Composite comp) {
                // Recursión para Composites anidados
                celdas.addAll(comp.celdasOcupadasActivas());
            } else if (c.isActivo()) {
                // Agregar píxeles activos
                celdas.add(new int[] { 
                    c.getRefX(), c.getRefY() 
                });
            }
        }
        return celdas;
    }
    
    @Override
    public void mover(int dx, int dy) {
        if (proyectil) {
            // Disparos: mover todos sin validar límites
            for (Component c : components) {
                c.mover(dx, dy);
            }
            return;
        }
        
        // NAVES: Validar límites ANTES de mover
        for (Component c : components) {
            int newX = c.getRefX() + dx;
            int newY = c.getRefY() + dy;
            if (newX < 0 || newX >= 100 || 
                newY < 0 || newY >= 60) {
                return;  // No permitir movimiento fuera
            }
        }
        
        // Guardar posiciones antiguas
        int[] oldPositionsX = 
            new int[components.size()];
        int[] oldPositionsY = 
            new int[components.size()];
        for (int i = 0; i < components.size(); i++) {
            oldPositionsX[i] = 
                components.get(i).getRefX();
            oldPositionsY[i] = 
                components.get(i).getRefY();
        }
        
        // Mover todos los componentes
        for (Component c : components) {
            c.mover(dx, dy);
        }
        
        // Notificar movimiento en bloque
        Espacio espacio = Espacio.getEspacio();
        espacio.notificarMovimientoJugadorCompleto(
            oldPositionsX, oldPositionsY, components);
    }
    
    // Referencia: esquina superior izquierda del bounding box
    @Override
    public int getRefX() {
        int min = Integer.MAX_VALUE;
        for (Component c : components) {
            min = Math.min(min, c.getRefX());
        }
        return min == Integer.MAX_VALUE ? 0 : min;
    }
    
    @Override
    public int getRefY() {
        int min = Integer.MAX_VALUE;
        for (Component c : components) {
            min = Math.min(min, c.getRefY());
        }
        return min == Integer.MAX_VALUE ? 0 : min;
    }
    
    @Override
    public boolean isActivo() {
        // Un Composite es activo si tiene al menos 
        // un componente activo
        return components.stream()
            .anyMatch(Component::isActivo);
    }
    
    @Override
    public void setActivo(boolean b) {
        // Propagar estado a todos los componentes
        for (Component c : components) {
            c.setActivo(b);
        }
    }
}
\end{lstlisting}

\textbf{Características principales:}

\begin{itemize}
    \item \textbf{Dos modos:} \texttt{proyectil=false} para naves, \texttt{proyectil=true} para disparos.
    
    \item \textbf{Movimiento recursivo:} El método \texttt{mover()} delega a todos sus componentes, formando una cascada.
    
    \item \textbf{Validación inteligente (naves):} Verifica que todos los píxeles quepan dentro del tablero (0-99 en X, 0-59 en Y) antes de permitir movimiento.
    
    \item \textbf{Movimiento sin restricción (disparos):} Los disparos se mueven libremente y se desactivan en sus píxeles cuando salen.
    
    \item \textbf{Notificación en bloque:} Naves notifican el movimiento de toda la estructura juntos, no píxel a píxel, para eficiencia.
    
    \item \textbf{Cálculo de referencias:} \texttt{getRefX/Y()} retornan la esquina superior izquierda, permitiendo posición de la forma completa.
    
    \item \textbf{Método \texttt{celdasOcupadasActivas()}:} Retorna todas las celdas ocupadas, esencial para colisiones.
\end{itemize}

\subsection{Construcción de Disparos Complejos}

La clase \texttt{Disparo} utiliza Composite para crear formas complejas:

\begin{lstlisting}
private Component construirCuerpo(int x, int y, String tipo) {
    switch (tipo) {
    case "flecha": {
        // Crear Composite para la flecha
        Composite comp = new Composite(true);
        comp.addComponent(new Pixel(x, y, true));
        comp.addComponent(new Pixel(x - 1, y + 1, true));
        comp.addComponent(new Pixel(x + 1, y + 1, true));
        return comp;  // Retorna como Component
    }
    case "rombo": {
        // Crear Composite para el rombo
        Composite comp = new Composite(true);
        comp.addComponent(new Pixel(x, y, true));
        comp.addComponent(new Pixel(x - 1, y + 1, true));
        comp.addComponent(new Pixel(x, y + 1, true));
        comp.addComponent(new Pixel(x + 1, y + 1, true));
        comp.addComponent(new Pixel(x, y + 2, true));
        return comp;
    }
    default:
        // Pixel simple
        return new Pixel(x, y, true);
    }
}
\end{lstlisting}

\textbf{Formas resultantes:}

\begin{center}
    \textbf{Píxel:} 1 celda\\
    $\bullet$\\[0.2cm]
    
    \textbf{Flecha:} 3 celdas en V\\
    $\bullet$\\
    $\bullet$ $\quad$ $\bullet$\\[0.2cm]
    
    \textbf{Rombo:} 5 celdas en rombo\\
    $\phantom{.}\bullet$\\
    $\bullet$ $\quad$ $\bullet$\\
    $\phantom{.}\bullet$\\
\end{center}

\subsection{Construcción de Naves}

Las naves usan un \texttt{Composite} como raíz para su estructura jerárquica:

\begin{lstlisting}
public class Nave1 extends Naves {
    private DisparoPixel estrategiaPixel = 
        new DisparoPixel();
    private DisparoFlecha estrategiaFlecha = 
        new DisparoFlecha();
    
    public Nave1(int x, int y, int velocidad) {
        super(x, y, velocidad);
        inicializarNaveJugador();
    }
    
    @Override
    public void construir() {
        int bx = x;
        int by = y;
        // Fila superior: 2 píxeles
        anadirComponente(new Pixel(bx, by));
        anadirComponente(new Pixel(bx + 2, by));
        // Fila inferior: 3 píxeles (base de T)
        anadirComponente(new Pixel(bx, by + 1));
        anadirComponente(new Pixel(bx + 1, by + 1));
        anadirComponente(new Pixel(bx + 2, by + 1));
    }
}
\end{lstlisting}

\textbf{Estructura de Nave1 (T invertida):}

\begin{lstlisting}
Diagrama visual en el tablero:
    
    $\bullet$ $\phantom{.}$ $\bullet$
    
    $\bullet$ $\bullet$ $\bullet$
\end{lstlisting}

\subsection{Diagrama de Clases Completo}

\begin{center}
    \includegraphics[width=1\textwidth]{images/composite_detailed.png}
    
    \textit{Figura 4: Estructura detallada de clases mostrando Composite (Component), Pixel, Disparo y sus relaciones.}
\end{center}

\subsection{Flujo de Ejecución: Movimiento de Nave}

El proceso de mover una nave demuestra la elegancia del Composite:

\begin{enumerate}
    \item \textbf{Usuario presiona tecla de movimiento} (izquierda o derecha)
    
    \item \textbf{Sistema solicita movimiento a nave:}
    \begin{lstlisting}
nave.mover(-1, 0);  // Izquierda
    \end{lstlisting}
    
    \item \textbf{Nave (Composite) verifica límites:}
    \begin{lstlisting}
for (Component c : components) {
    if (newX < 0 || newX >= 100) return;
}
    \end{lstlisting}
    
    \item \textbf{Naveve délega a todos sus componentes:}
    \begin{lstlisting}
for (Component c : components) {
    c.mover(dx, dy);  // Cada Pixel actualiza (x, y)
}
    \end{lstlisting}
    
    \item \textbf{Cada Pixel se mueve:} \texttt{x--}
    
    \item \textbf{Tablero se actualiza:} Los Pixels notifican movimiento para redibujarse.
    
    \item \textbf{Próximo frame:} Colisiones se detectan en nuevas posiciones.
\end{enumerate}

\textbf{Diagrama de secuencia:}

\begin{verbatim}
Usuario (flecha izquierda)
   |
   v
Nave.mover(-1, 0)
   |
   |-- Validar límites
   |
   v
Composite.mover(-1, 0)
   |
   |-- Pixel(0, 0).mover(-1, 0)  --> x = -1
   |-- Pixel(2, 0).mover(-1, 0)  --> x = 1
   |-- Pixel(0, 1).mover(-1, 0)  --> x = -1
   |-- Pixel(1, 1).mover(-1, 0)  --> x = 0
   |-- Pixel(2, 1).mover(-1, 0)  --> x = 1
   |
   v
Notificar movimiento completo (5 píxeles)
   |
   v
Tablero renderiza nuevas posiciones
\end{verbatim}

\subsection{Flujo de Ejecución: Colisión}

La detección de colisiones aprovecha el Composite:

\begin{enumerate}
    \item \textbf{Cada frame, sistema consulta celdas ocupadas:}
    \begin{lstlisting}
List<int[]> celdasNavesActivas = 
    nave.getCeldasOcupadas();
    \end{lstlisting}
    
    \item \textbf{Sistema consulta celdas de disparos enemigos:}
    \begin{lstlisting}
List<int[]> celdasDisparosEnemigos = 
    enemigo.getDisparos()
        .stream()
        .flatMap(d -> d.getCeldas().stream())
        .collect(toList());
    \end{lstlisting}
    
    \item \textbf{Comparar para encontrar intersecciones:}
    \begin{lstlisting}
for (int[] celdaNave : celdasNavesActivas) {
    for (int[] celdaDisparo : 
            celdasDisparosEnemigos) {
        if (Arrays.equals(celdaNave, celdaDisparo)) {
            // COLISIÓN DETECTADA
            nave.perderVida();
            disparo.desactivar();
        }
    }
}
    \end{lstlisting}
\end{enumerate}

\subsection{Ventajas del Patrón Composite}

\begin{enumerate}
    \item \textbf{Uniformidad extrema:} Píxeles individuales y estructuras complejas usan la misma interfaz.
    
    \item \textbf{Operaciones transversales simples:} Mover una nave mueve automáticamente todos sus píxeles sin bucles adicionales.
    
    \item \textbf{Detección de colisiones unificada:} Un método \texttt{celdasOcupadasActivas()} funciona para cualquier estructura.
    
    \item \textbf{Flexibilidad geométrica:} Crear nuevas formas solo requiere construir diferentes Composites en \texttt{construir()}.
    
    \item \textbf{Eficiencia en notificaciones:} Naves notifican bloques de píxeles, no individualmente.
    
    \item \textbf{Gestión de ciclo de vida automática:} Píxeles de disparo se desactivan automáticamente al salir del tablero.
    
    \item \textbf{Reutilización masiva:} El patrón sirve para naves, disparos complejos y futuros elementos (enemigos complejos).
\end{enumerate}

\subsection{Comparativa: Sin vs Con Composite}

\textbf{Sin Composite (Acoplado y repetitivo):}

\begin{lstlisting}
public class NaveAntiguA {
    private int[] pixelX = {0, 2, 0, 1, 2};
    private int[] pixelY = {0, 0, 1, 1, 1};
    
    public void mover(int dx, int dy) {
        // Validación manual
        for (int i = 0; i < pixelX.length; i++) {
            int newX = pixelX[i] + dx;
            if (newX < 0 || newX >= 100) return;
        }
        // Movimiento manual
        for (int i = 0; i < pixelX.length; i++) {
            pixelX[i] += dx;
            pixelY[i] += dy;
        }
    }
    
    public List<int[]> getCeldas() {
        List<int[]> celdas = new ArrayList<>();
        for (int i = 0; i < pixelX.length; i++) {
            celdas.add(new int[] {pixelX[i], pixelY[i]});
        }
        return celdas;
    }
    // Repetir esto para cada tipo de nave...
}
\end{lstlisting}

\textbf{Con Composite (Elegante y escalable):}

\begin{lstlisting}
public class Nave1 extends Naves {
    public void construir() {
        anadirComponente(new Pixel(0, 0));
        anadirComponente(new Pixel(2, 0));
        anadirComponente(new Pixel(0, 1));
        anadirComponente(new Pixel(1, 1));
        anadirComponente(new Pixel(2, 1));
    }
}

// Todas las operaciones heredadas:
nave.mover(-1, 0);           // Funciona
nave.getCeldasOcupadas();    // Funciona
nave.getCeldasActivas();     // Funciona
// Sin duplicar código
\end{lstlisting}

\newpage
\section{Integración: Cómo Trabajan Juntos los Patrones}

\subsection{Flujo Completo de Juego}

Los tres patrones trabajan en conjunto para formar un sistema cohesivo:

\begin{enumerate}
    \item \textbf{Inicio del juego:}
    \begin{enumerate}
        \item \textbf{Factory} crea la nave seleccionada mediante \texttt{NaveFactory.generate()}.
        \item La nave contiene un \textbf{Composite} como estructura.
        \item El jugador tiene acceso a las \textbf{Strategies} de disparo.
    \end{enumerate}
    
    \item \textbf{Durante el juego - Movimiento:}
    \begin{enumerate}
        \item Usuario presiona tecla.
        \item Nave (Composite) invoca \texttt{mover()}.
        \item Todos los píxeles se desplazan, gracua al \textbf{Composite}.
    \end{enumerate}
    
    \item \textbf{Durante el juego - Disparo:}
    \begin{enumerate}
        \item Usuario selecciona tipo de disparo (Strategy).
        \item Sistema verifica munición en la estrategia.
        \item Crea Disparo con Composite de forma según la Strategy.
        \item Disparo sube usando movimiento \textbf{Composite}.
    \end{enumerate}
    
    \item \textbf{Detección de colisiones:}
    \begin{enumerate}
        \item Sistema obtiene todas las celdas ocupadas (Composite).
        \item Compara celdas de naves vs disparos vs enemigos.
        \item Resuelve colisiones.
    \end{enumerate}
\end{enumerate}

\subsection{Diagrama de Interacción}

\begin{verbatim}
INICIO
   |
   v
Factory.generate("Nave1")
   |
   +---> new Nave1(50, 55, 1)
   |         |
   |         v
   |     Composite raíz
   |     |
   |     +---> Pixel(0,0)
   |     +---> Pixel(2,0)
   |     +---> Pixel(0,1)
   |     +---> Pixel(1,1)
   |     +---> Pixel(2,1)
   |
   v
Nave lista con estrategias
   |
   +---> DisparoPixel       (Strategy)
   +---> DisparoFlecha      (Strategy)
   
DURANTE JUEGO
   |
   +-- Usuario mueve <---> nave.mover() <---> Composite
   |
   +-- Usuario dispara:
   |       |
   |       v
   |   StrategyDisparo.tieneMunicion()?
   |       |
   |       v
   |   Disparo nuevo
   |       |
   |       v
   |   new Pixel/Composite según tipo
   |       |
   |       v
   |   disparos.add(nuevoDisparo)
   |
   +-- Colisiones:
   |       |
   |       v
   |   nave.getCeldasOcupadas()
   |   disparo.getCeldasOcupadas()
   |   enemigo.getCeldasOcupadas()
   |       |
   |       v
   |   Comparar e resolver
\end{verbatim}

\section{Conclusión}

La combinación de los patrones \textbf{Factory}, \textbf{Strategy} y \textbf{Composite} en \textit{Space Invaders} demuestra una aplicación sólida de los principios de diseño orientado a objetos:

\begin{enumerate}
    \item \textbf{Factory:} Gestiona la creación centralizada de naves, permitiendo agregar nuevas variantes sin modificar código existente.
    
    \item \textbf{Strategy:} Permite cambiar comportamientos de disparo en tiempo de ejecución, cada estrategia manejando su propio estado.
    
    \item \textbf{Composite:} Facilita el manejo uniforme de estructuras jerárquicas, simplificando movimiento, colisiones y renderizado.
\end{enumerate}

\subsection{Principios SOLID Aplicados}

\begin{itemize}
    \item \textbf{Single Responsibility:} Cada clase tiene una responsabilidad clara.
    \item \textbf{Open/Closed:} El sistema es abierto para extensión (nuevas naves, estrategias) pero cerrado para modificación.
    \item \textbf{Liskov Substitution:} Las implementaciones de Strategy y Component son intercambiables.
    \item \textbf{Interface Segregation:} Las interfaces son específicas y no fuerzan métodos innecesarios.
    \item \textbf{Dependency Inversion:} El código depende de abstracciones (interfaces), no de implementaciones.
\end{itemize}

\subsection{Beneficios Finales}

\begin{itemize}
    \item \textbf{Mantenibilidad:} Cambios localizados, bajo acoplamiento.
    \item \textbf{Extensibilidad:} Nuevas características sin afectar código existente.
    \item \textbf{Clarity:} Código autoexplicativo, patrones conocidos.
    \item \textbf{Robustez:} Estructura clara reduce errores.
    \item \textbf{Testabilidad:} Componentes aislados, fáciles de probar.
\end{itemize}

Esta es una implementación didáctica y profesional de cómo tres patrones fundamentales resuelven problemas reales en el desarrollo de videos juegos, demostrando que una buena arquitectura es la base del software mantenible.

\end{document}
    \item Un disparo complejo puede ser una estructura compuesta de varios píxeles.
\end{itemize}

Gracias a este diseño, las operaciones como mover, activar o detectar colisiones se pueden aplicar de manera uniforme a cualquier elemento, independientemente de su complejidad.

\subsection{Ventajas}

El uso del patrón Composite aporta múltiples beneficios:

\begin{itemize}
    \item Permite tratar estructuras simples y complejas de forma uniforme.
    \item Facilita la implementación de operaciones recursivas como el movimiento.
    \item Simplifica la detección de colisiones al unificar el acceso a las celdas ocupadas.
    \item Favorece la reutilización de código.
    \item Permite crear estructuras geométricas complejas de forma flexible.
\end{itemize}

\section{Conclusión}

La combinación de los patrones Factory, Strategy y Composite en \textit{Space Invaders} demuestra una aplicación sólida de los principios de diseño orientado a objetos.

Cada patrón resuelve un problema específico:

\begin{itemize}
    \item Factory gestiona la creación de objetos de forma centralizada.
    \item Strategy permite variar comportamientos dinámicamente.
    \item Composite facilita el manejo de estructuras jerárquicas complejas.
\end{itemize}

En conjunto, estos patrones contribuyen a un sistema más modular, extensible y fácil de mantener, alineado con principios fundamentales como el bajo acoplamiento y la alta cohesión.

\end{document}